\chapter{The natural numbers}

Everything starts with the natural numbers. Since \texttt{lean} already has a type of natural
numbers, we define our own copy, called $\MyNat$, from scratch. To make sure we do not secretly use
results that are already in mathlib, the file defining $\MyNat$ does not import anything.

\section{Definition}

\begin{definition}
    \label{MyNat}
    \lean{MyNat}
    \leanok
    We define $\MyNat$ as the inductive type with two constructors:
    \begin{itemize}
        \item a term $0$ (called \texttt{zero});
        \item for every $n \in \MyNat$, a term $\mathrm{succ}(n)$ (the \emph{successor} of $n$).
    \end{itemize}
    In other words, every natural number is obtained from $0$ by applying the successor function a
    finite number of times. We write $1$ for $\mathrm{succ}(0)$.
\end{definition}

\begin{definition}
    \label{MyNat.add}
    \lean{MyNat.add}
    \uses{MyNat}
    \leanok
    We define \emph{addition} on $\MyNat$ by recursion on the second argument:
    \[
    a + 0 = a, \qquad a + \mathrm{succ}(b) = \mathrm{succ}(a + b).
    \]
\end{definition}

\begin{definition}
    \label{MyNat.mul}
    \lean{MyNat.mul}
    \uses{MyNat.add}
    \leanok
    We define \emph{multiplication} on $\MyNat$ by recursion on the second argument:
    \[
    a * 0 = 0, \qquad a * \mathrm{succ}(b) = a * b + a.
    \]
\end{definition}

\section{Commutative semiring structure}

All the lemmas of this section are proved by induction, using the recursive definitions of addition
and multiplication. Note that $a + 0 = a$ and $a * 0 = 0$ hold \emph{by definition}, while the
symmetric statements $0 + a = a$ and $0 * a = 0$ need a proof.

\begin{lemma}
    \label{MyNat.zero_add}
    \lean{MyNat.zero_add}
    \leanok
    For all $a$ in $\MyNat$ we have $0 + a = a$.
\end{lemma}
\begin{proof}
    \uses{MyNat.add}
    \leanok
    By induction on $a$.
\end{proof}

\begin{lemma}
    \label{MyNat.add_assoc}
    \lean{MyNat.add_assoc}
    \leanok
    Addition on $\MyNat$ is associative.
\end{lemma}
\begin{proof}
    \uses{MyNat.add}
    \leanok
    By induction on the last variable.
\end{proof}

\begin{lemma}
    \label{MyNat.add_comm}
    \lean{MyNat.add_comm}
    \leanok
    Addition on $\MyNat$ is commutative.
\end{lemma}
\begin{proof}
    \uses{MyNat.add, MyNat.zero_add}
    \leanok
    By induction.
\end{proof}

\begin{lemma}
    \label{MyNat.left_distrib}
    \lean{MyNat.left_distrib}
    \leanok
    For all $a$, $b$ and $c$ in $\MyNat$ we have $a * (b + c) = a * b + a * c$.
\end{lemma}
\begin{proof}
    \uses{MyNat.add, MyNat.mul, MyNat.add_assoc, MyNat.add_comm}
    \leanok
    By induction on $c$.
\end{proof}

\begin{lemma}
    \label{MyNat.right_distrib}
    \lean{MyNat.right_distrib}
    \leanok
    For all $a$, $b$ and $c$ in $\MyNat$ we have $(a + b) * c = a * c + b * c$.
\end{lemma}
\begin{proof}
    \uses{MyNat.add, MyNat.mul, MyNat.add_assoc, MyNat.add_comm}
    \leanok
    By induction on $c$.
\end{proof}

\begin{lemma}
    \label{MyNat.mul_assoc}
    \lean{MyNat.mul_assoc}
    \leanok
    Multiplication on $\MyNat$ is associative.
\end{lemma}
\begin{proof}
    \uses{MyNat.mul, MyNat.left_distrib}
    \leanok
    By induction, using left distributivity.
\end{proof}

\begin{lemma}
    \label{MyNat.mul_comm}
    \lean{MyNat.mul_comm}
    \leanok
    Multiplication on $\MyNat$ is commutative.
\end{lemma}
\begin{proof}
    \uses{MyNat.mul}
    \leanok
    By induction.
\end{proof}

\begin{lemma}
    \label{MyNat.one_mul}
    \lean{MyNat.one_mul}
    \leanok
    For all $a$ in $\MyNat$ we have $1 * a = a$.
\end{lemma}
\begin{proof}
    \uses{MyNat.mul}
    \leanok
    By induction, recalling that $1 = \mathrm{succ}(0)$.
\end{proof}

\begin{proposition}
    \label{MyNat.instCommSemiring}
    \lean{MyNat.instCommSemiring}
    \leanok
    $\MyNat$ with addition and multiplication is a commutative semiring.
\end{proposition}
\begin{proof}
    \uses{MyNat.add, MyNat.mul, MyNat.zero_add, MyNat.add_assoc, MyNat.add_comm,
    MyNat.left_distrib, MyNat.right_distrib, MyNat.mul_assoc, MyNat.mul_comm, MyNat.one_mul}
    \leanok
    We collect all the lemmas above: addition is associative and commutative with neutral element
    $0$, multiplication is associative and commutative with neutral element $1$, multiplication
    distributes over addition and $0$ is absorbing for multiplication.
\end{proof}

\begin{lemma}
    \label{MyNat.zero_ne_one}
    \lean{MyNat.zero_ne_one}
    \leanok
    In $\MyNat$ we have $0 \neq 1$.
\end{lemma}
\begin{proof}
    \uses{MyNat}
    \leanok
    The constructors $\mathrm{zero}$ and $\mathrm{succ}$ are distinct, and $1 = \mathrm{succ}(0)$.
\end{proof}

\begin{lemma}
    \label{MyNat.mul_ne_zero}
    \lean{MyNat.mul_ne_zero}
    \leanok
    Let $a$ and $b$ in $\MyNat$ be such that $a \neq 0$ and $b \neq 0$. Then $a * b \neq 0$.
\end{lemma}
\begin{proof}
    \uses{MyNat.mul}
    \leanok
    If $a \neq 0$ and $b \neq 0$ we can write $a = \mathrm{succ}(a')$ and $b = \mathrm{succ}(b')$, and
    then $a * b = \mathrm{succ}(a' * b' + a' + b')$ is a successor, hence nonzero.
\end{proof}

\section{The order}

\begin{definition}
    \label{MyNat.le}
    \lean{MyNat.le}
    \uses{MyNat.add}
    \leanok
    Let $a$ and $b$ in $\MyNat$. We write $a \leq b$ if there exists a natural number $x$ such that
    \[
    b = a + x.
    \]
\end{definition}

\begin{lemma}
    \label{MyNat.le_refl}
    \lean{MyNat.le_refl}
    \leanok
    The relation $\leq$ on $\MyNat$ is reflexive.
\end{lemma}
\begin{proof}
    \uses{MyNat.le}
    \leanok
    We can take $x = 0$.
\end{proof}

\begin{lemma}
    \label{MyNat.le_trans}
    \lean{MyNat.le_trans}
    \leanok
    The relation $\leq$ on $\MyNat$ is transitive.
\end{lemma}
\begin{proof}
    \uses{MyNat.le, MyNat.add_assoc}
    \leanok
    If $b = a + x$ and $c = b + y$ then $c = a + (x + y)$.
\end{proof}

\begin{lemma}
    \label{MyNat.le_antisymm}
    \lean{MyNat.le_antisymm}
    \leanok
    The relation $\leq$ on $\MyNat$ is antisymmetric.
\end{lemma}
\begin{proof}
    \uses{MyNat.le}
    \leanok
    If $b = a + x$ and $a = b + y$ then $a = a + (x + y)$, so $x + y = 0$ and hence $x = y = 0$,
    giving $a = b$.
\end{proof}

\begin{lemma}
    \label{MyNat.le_total}
    \lean{MyNat.le_total}
    \leanok
    The order $\leq$ on $\MyNat$ is total.
\end{lemma}
\begin{proof}
    \uses{MyNat.le}
    \leanok
    By induction one shows that for all $a$ and $b$ either $a \leq b$ or $b \leq a$.
\end{proof}

\begin{lemma}
    \label{MyNat.linearOrder}
    \lean{MyNat.linearOrder}
    \leanok
    We have that $\MyNat$ with $\leq$ is a \emph{linear order}.
\end{lemma}
\begin{proof}
    \uses{MyNat.le_refl, MyNat.le_trans, MyNat.le_antisymm, MyNat.le_total}
    \leanok
    Clear from the lemmas above.
\end{proof}

\section{Interaction between the order and the algebraic structure}

\begin{lemma}
    \label{MyNat.add_le_add_left}
    \lean{MyNat.add_le_add_left}
    \leanok
    Let $a$, $b$ and $c$ in $\MyNat$ be such that $a \leq b$. Then $a + c \leq b + c$.
\end{lemma}
\begin{proof}
    \uses{MyNat.le, MyNat.add}
    \leanok
    Let $x$ be such that $b = a + x$. The same $x$ shows that $a + c \leq b + c$.
\end{proof}

\begin{lemma}
    \label{MyNat.mul_lt_mul_of_pos_left}
    \lean{MyNat.mul_lt_mul_of_pos_left}
    \leanok
    Let $a$, $b$ and $c$ in $\MyNat$ be such that $0 < c$ and $a < b$. Then $c * a < c * b$.
\end{lemma}
\begin{proof}
    \uses{MyNat.le, MyNat.mul, MyNat.left_distrib, MyNat.mul_ne_zero}
    \leanok
    Write $b = a + x$ with $x \neq 0$. Then $c * b = c * a + c * x$ with $c * x \neq 0$, so
    $c * a < c * b$.
\end{proof}

\begin{lemma}
    \label{MyNat.mul_lt_mul_of_pos_right}
    \lean{MyNat.mul_lt_mul_of_pos_right}
    \leanok
    Let $a$, $b$ and $c$ in $\MyNat$ be such that $0 < c$ and $a < b$. Then $a * c < b * c$.
\end{lemma}
\begin{proof}
    \uses{MyNat.mul_lt_mul_of_pos_left, MyNat.mul_comm}
    \leanok
    Identical to the previous lemma using commutativity of multiplication.
\end{proof}
